\documentclass[12pt,a4paper]{article}
\usepackage[utf8]{inputenc}
\usepackage[spanish]{babel}
\usepackage[T1]{fontenc}
\usepackage{geometry}
\geometry{top=2.5cm, bottom=2.5cm, left=2.5cm, right=2.5cm}
\usepackage{enumitem}
\usepackage{amssymb}
\usepackage{amsmath}
\usepackage{booktabs}
\usepackage{fancyhdr}
\usepackage{listings}
\usepackage{xcolor}
\usepackage{array}
\usepackage{tabularx}
\usepackage{mdframed}
\usepackage{tcolorbox}
\usepackage{hyperref}
\hypersetup{colorlinks=false}

% -------------------------------------------------------
% Colores y estilo de código Python
% -------------------------------------------------------
\definecolor{codegreen}{rgb}{0.1,0.5,0.1}
\definecolor{codegray}{rgb}{0.4,0.4,0.4}
\definecolor{codepurple}{rgb}{0.5,0.0,0.5}
\definecolor{backcolour}{rgb}{0.97,0.97,0.97}
\definecolor{anahuacblue}{RGB}{0,51,102}

\lstdefinestyle{PythonStyle}{
    backgroundcolor=\color{backcolour},
    commentstyle=\color{codegreen},
    keywordstyle=\color{codepurple}\bfseries,
    numberstyle=\tiny\color{codegray},
    stringstyle=\color{codegreen},
    basicstyle=\ttfamily\small,
    breakatwhitespace=false,
    breaklines=true,
    keepspaces=true,
    numbers=left,
    numbersep=5pt,
    showspaces=false,
    showstringspaces=false,
    showtabs=false,
    tabsize=4,
    language=Python
}

% -------------------------------------------------------
% Encabezado y pie de página
% -------------------------------------------------------
\pagestyle{fancy}
\fancyhf{}
\lhead{\small\textbf{Algoritmos y Programación} --- Examen Parcial2}
\rhead{\small\textbf{Universidad Anáhuac México}}
\lfoot{\small Prof.\ Dr.\ BARSEKH-ONJI Aboud}
\rfoot{\small Página \thepage}
\renewcommand{\headrulewidth}{0.6pt}
\renewcommand{\footrulewidth}{0.4pt}

% -------------------------------------------------------
% Comando para líneas de respuesta
% -------------------------------------------------------
\newcommand{\answerline}{\vspace{0.5cm}\noindent\rule{\linewidth}{0.3pt}\vspace{0.3cm}}

% -------------------------------------------------------
\begin{document}

% ===========================
%  ENCABEZADO INSTITUCIONAL
% ===========================
\thispagestyle{fancy}

\begin{center}
  {\color{anahuacblue}\rule{\linewidth}{2pt}}\\[0.3cm]
  {\Large\bfseries UNIVERSIDAD ANÁHUAC MÉXICO}\\[0.1cm]
  {\large Facultad de Ingeniería}\\[0.1cm]
  {\large\bfseries Algoritmos y Programación}\\[0.2cm]
  {\LARGE\bfseries EXAMEN}\\[0.1cm]
  {\large Prof.\ Dr.\ BARSEKH-ONJI Aboud}\\[0.2cm]
  {\color{anahuacblue}\rule{\linewidth}{2pt}}
\end{center}

\vspace{0.4cm}

% Datos del alumno
\noindent
\begin{tabularx}{\linewidth}{X X}
  \textbf{Nombre del alumno:} \hrulefill & \textbf{Matrícula:} \hrulefill \\[0.6cm]
  \textbf{Fecha:} \hrulefill & \textbf{Grupo:} \hrulefill \\[0.3cm]
\end{tabularx}

\vspace{0.3cm}

% ===========================
%  INSTRUCCIONES GENERALES
% ===========================
\begin{tcolorbox}[colback=anahuacblue!8, colframe=anahuacblue, title={\bfseries Instrucciones Generales}, fonttitle=\bfseries\color{white}]
\begin{itemize}[leftmargin=1.5em, itemsep=2pt]
  \item Duración del examen: \textbf{90 minutos}.
  \item \textbf{NO} se permite el uso de computadora, laptop ni teléfono inteligente de ningún tipo.
  \item Se permite el uso de \textbf{calculadora} si es necesario.
  \item En la \textbf{Sección A} (opción múltiple), una pregunta puede tener \textbf{una o varias respuestas correctas}. \\ Debes marcar \textbf{TODAS} las opciones correctas para obtener el puntaje completo de la pregunta.
  \item En la \textbf{Sección B} (problemas), escribe tu código con letra \textbf{clara y legible}. El código ilegible no será calificado.
  \item Respeta la \textbf{indentación} de Python al escribir tu código (usa 4 espacios por nivel).
\end{itemize}
\end{tcolorbox}

\vspace{0.4cm}

% ===========================
%  TABLA DE CALIFICACIÓN
% ===========================
\begin{center}
\begin{tabular}{lcc}
  \toprule
  \textbf{Sección} & \textbf{Puntos posibles} & \textbf{Puntos obtenidos} \\
  \midrule
  A — Opción múltiple (8 preguntas) & 40 & \\
  B — Problema 1 & 30 & \\
  B — Problema 2 & 30 & \\
  \midrule
  \textbf{Total} & \textbf{100} & \\
  \bottomrule
\end{tabular}
\end{center}

\newpage

% ===========================
%  SECCIÓN A — OPCIÓN MÚLTIPLE
% ===========================
\section*{\color{anahuacblue}Sección A — Opción Múltiple \normalsize(5 puntos cada pregunta. Puede haber varias respuestas correctas. Marca \textbf{TODAS} las correctas.)}

\medskip

% --- Pregunta 1 ---
\noindent\textbf{Pregunta 1.} ¿Cuál(es) de los siguientes fragmentos de código imprime los números \texttt{1, 3, 5, 7, 9}?

\begin{enumerate}[label=\alph*), leftmargin=2em, itemsep=4pt]
  \item[$\square$ a)] \lstinline|for i in range(1, 10, 2): print(i)|
  \item[$\square$ b)] \lstinline|for i in range(0, 10, 2): print(i)|
  \item[$\square$ c)] \lstinline|for i in range(1, 11, 2): print(i)|
  \item[$\square$ d)] \lstinline|for i in range(9, 0, -2): print(i)|
\end{enumerate}

\bigskip

% --- Pregunta 2 ---
\noindent\textbf{Pregunta 2.} Analiza el siguiente programa:
\begin{lstlisting}[style=PythonStyle]
contador = 0
valores = [3, -1, 7, -4, 5, -2, 8]
for v in valores:
    if v > 0:
        contador += 1
print(contador)
\end{lstlisting}
¿Qué valor imprime el programa?

\begin{enumerate}[label=\alph*), leftmargin=2em, itemsep=4pt]
  \item[$\square$ a)] 3
  \item[$\square$ b)] 4
  \item[$\square$ c)] 7
  \item[$\square$ d)] 18
\end{enumerate}

\bigskip

% --- Pregunta 3 ---
\noindent\textbf{Pregunta 3.} Dada la lista \lstinline|datos = [10, 3, 8, 3, 10]|, ¿cuál(es) de las siguientes afirmaciones son correctas?

\begin{enumerate}[label=\alph*), leftmargin=2em, itemsep=4pt]
  \item[$\square$ a)] \lstinline|len(datos)| devuelve \texttt{5}
  \item[$\square$ b)] \lstinline|sum(datos)| devuelve \texttt{34}
  \item[$\square$ c)] \lstinline|max(datos)| devuelve \texttt{10}
  \item[$\square$ d)] \lstinline|datos[4]| devuelve \texttt{3}
\end{enumerate}

\bigskip

% --- Pregunta 4 ---
\noindent\textbf{Pregunta 4.} Sea \lstinline|lista = [5, 10, 15, 20]|. Después de ejecutar:
\begin{lstlisting}[style=PythonStyle]
elemento = lista.pop(1)
lista.append(elemento * 2)
\end{lstlisting}
¿Cuál es el estado final de \lstinline|lista|?

\begin{enumerate}[label=\alph*), leftmargin=2em, itemsep=4pt]
  \item[$\square$ a)] \lstinline|[5, 10, 15, 20, 20]|
  \item[$\square$ b)] \lstinline|[5, 15, 20, 20]|
  \item[$\square$ c)] \lstinline|[5, 15, 20, 10]|
  \item[$\square$ d)] \lstinline|[10, 5, 15, 20]|
\end{enumerate}

\bigskip

% --- Pregunta 5 ---
\noindent\textbf{Pregunta 5.} Analiza el siguiente bloque condicional:
\begin{lstlisting}[style=PythonStyle]
puntaje = 85
if puntaje >= 90:
    print("A")
elif puntaje >= 80:
    print("B")
elif puntaje >= 70:
    print("C")
else:
    print("F")
\end{lstlisting}
¿Qué imprime el programa? ¿Y cuál(es) de las siguientes afirmaciones también son verdaderas?

\begin{enumerate}[label=\alph*), leftmargin=2em, itemsep=4pt]
  \item[$\square$ a)] Imprime \texttt{B}
  \item[$\square$ b)] El bloque \lstinline|elif puntaje >= 70| nunca se evalúa en esta ejecución
  \item[$\square$ c)] Si \texttt{puntaje} fuera \texttt{90}, imprimiría tanto \texttt{A} como \texttt{B}
  \item[$\square$ d)] Si \texttt{puntaje} fuera \texttt{65}, imprimiría \texttt{F}
\end{enumerate}

\bigskip

% --- Pregunta 6 ---
\noindent\textbf{Pregunta 6.} Respecto a \lstinline|random.choice(secuencia)|, ¿cuál(es) afirmaciones son correctas?

\begin{enumerate}[label=\alph*), leftmargin=2em, itemsep=4pt]
  \item[$\square$ a)] Devuelve un elemento aleatorio de la secuencia
  \item[$\square$ b)] \lstinline|random.choice([10, 20, 30])| puede devolver \texttt{20}
  \item[$\square$ c)] Requiere \lstinline|import random| antes de usarse
  \item[$\square$ d)] \lstinline|random.choice([])| devuelve \texttt{None} sin error
\end{enumerate}

\bigskip

% --- Pregunta 7 ---
\noindent\textbf{Pregunta 7.} Considera el siguiente ciclo \texttt{while}:
\begin{lstlisting}[style=PythonStyle]
acum = 0
n = 1
while n <= 5:
    acum += n
    n += 1
print(acum)
\end{lstlisting}
¿Qué imprime el programa? Selecciona también las afirmaciones correctas.

\begin{enumerate}[label=\alph*), leftmargin=2em, itemsep=4pt]
  \item[$\square$ a)] Imprime \texttt{15}
  \item[$\square$ b)] El ciclo se ejecuta exactamente 5 veces
  \item[$\square$ c)] Imprime \texttt{10}
  \item[$\square$ d)] Al terminar el ciclo, \texttt{n} tiene el valor \texttt{6}
\end{enumerate}

\bigskip

% --- Pregunta 8 ---
\noindent\textbf{Pregunta 8.} ¿Cuál(es) de los siguientes fragmentos construyen correctamente una lista con los cuadrados de los números del 1 al 5, es decir \lstinline|[1, 4, 9, 16, 25]|?

\vspace{4pt}
\noindent$\square$\ \textbf{a)}
\begin{lstlisting}[style=PythonStyle, numbers=none, xleftmargin=1.5em]
cuadrados = []
for i in range(1, 6):
    cuadrados.append(i ** 2)
\end{lstlisting}

\vspace{2pt}
\noindent$\square$\ \textbf{b)}
\begin{lstlisting}[style=PythonStyle, numbers=none, xleftmargin=1.5em]
cuadrados = []
i = 1
while i <= 5:
    cuadrados.append(i * i)
    i += 1
\end{lstlisting}

\vspace{2pt}
\noindent$\square$\ \textbf{c)}
\begin{lstlisting}[style=PythonStyle, numbers=none, xleftmargin=1.5em]
cuadrados = []
for i in range(6):
    cuadrados.append(i ** 2)
\end{lstlisting}

\vspace{2pt}
\noindent$\square$\ \textbf{d)}
\begin{lstlisting}[style=PythonStyle, numbers=none, xleftmargin=1.5em]
cuadrados = [1, 4, 9, 16, 25]
\end{lstlisting}

\newpage

% ===========================
%  SECCIÓN B — PROBLEMAS
% ===========================
\section*{\color{anahuacblue}Sección B — Problemas de Programación \normalsize(30 puntos cada problema)}

\medskip

% --- Problema 1 ---
\subsection*{Problema 1 — Procesamiento de calificaciones (30 puntos)}

\noindent Se tiene la siguiente lista de calificaciones de un grupo de alumnos:

\begin{lstlisting}[style=PythonStyle]
calificaciones = [78, 92, 55, 63, 88, 45, 71, 96, 60, 83]
\end{lstlisting}

\noindent Escribe un programa en Python completo que realice \textbf{todas} las siguientes tareas:

\begin{enumerate}[leftmargin=2em, itemsep=5pt]
  \item Usando un ciclo \texttt{while}, recorre la lista e imprime cada calificación junto con la etiqueta \texttt{'Aprobado'} si es mayor o igual a 70, o \texttt{'Reprobado'} si es menor a 70. Usa \texttt{if - else}.
  \item Calcule e imprima el \textbf{promedio} del grupo usando \lstinline|sum()| y \lstinline|len()|.
  \item Imprima la \textbf{calificación más alta} con \lstinline|max()|.
  \item Elimine la \textbf{última} calificación de la lista con \lstinline|pop()| e imprima qué calificación fue eliminada.
  \item Agregue la calificación \texttt{74} a la lista con \lstinline|append()| e imprima la lista final.
\end{enumerate}


% --- Problema 2 ---
\subsection*{Problema 2 — Simulación de dados (30 puntos)}

\noindent Escribe un programa en Python que simule el lanzamiento de un dado de 6 caras para un jugador y realice lo siguiente:

\begin{enumerate}[leftmargin=2em, itemsep=5pt]
  \item Crea una lista vacía llamada \lstinline|resultados|.
  \item Usando un ciclo \texttt{for} con \lstinline|range(10)|, genera 10 lanzamientos de dado usando \lstinline|random.randint(1, 6)| y agrega cada resultado a \lstinline|resultados| con \lstinline|append()|.
  \item Imprime la lista \lstinline|resultados|.
  \item Usando un ciclo \texttt{for} sobre \lstinline|resultados|, cuenta cuántos lanzamientos resultaron en \texttt{6} (usa \texttt{if}) e imprime ese conteo con el mensaje \texttt{"Número de seises: X"}.
  \item Usando \texttt{if / elif / else}, evalúa el \textbf{total acumulado} de los 10 lanzamientos (usa \lstinline|sum()|) e imprime:
    \begin{itemize}
      \item \texttt{"Puntuación alta"} si el total es mayor o igual a 40
      \item \texttt{"Puntuación media"} si el total está entre 25 y 39 (inclusive)
      \item \texttt{"Puntuación baja"} si el total es menor a 25
    \end{itemize}
\end{enumerate}


\vfill
\begin{center}
  {\color{anahuacblue}\rule{0.5\linewidth}{0.8pt}}\\[0.2cm]
  {\small\itshape ``La tecnología es una herramienta; el pensamiento lógico es el verdadero poder.''}\\
  {\small Prof.\ Dr.\ BARSEKH-ONJI Aboud --- Universidad Anáhuac México}
\end{center}

\end{document}
